
\input texsis
\paper

\overfullrule=0pt
\def\frac#1#2{#1\over #2}
%
\def\bmatrix#1{\left[ \matrix{#1} \right]}
\def\be{{\beta}}
\def\toone{{t+1}}

%\def\frac#1#2{#1\over #2}
%\magnification=1200
\overfullrule=0pt

\line{PhD Macroeconomics \hfil NYU}
\line{November 24, 2020 \hfil Thomas Sargent}

\centerline{\bf  More Fun Problems}
%\centerline{\bf Homework 3}


\medskip

\noindent{\bf Instructions:}   These problems explore ideas we have studied.


\medskip

\noindent{\bf 0.}  Where $c_t$ is consumption at date $t$ and $b_t$ is one-period debt owed at time $t$,
a consumer chooses a consumption, borrowing plan $\{c_t, b_{t+1}\}_{t=0}^T$
to maximize
$$ \sum_{t=0}^T\beta^t \ln(c_t),  \ \  \beta \in (0,1) $$
subject to the budget constraints
$$  c_t + b_t \leq y_t + \beta b_{t+1} , \ \ t = 0, \ldots, T $$
where $y_t = \lambda^t , \ \lambda \in (0,1)$, initial one-period risk-free debt $b_0 = 0$, and we impose
 $b_{T+1} \leq 0$.   


\medskip
\noindent{\bf a.} Verify that it is optimal for the consumer to set $b_{T+1} =0$.

\medskip
\noindent{\bf b.} Argue that  because the consumer can afford to set $c_t >0$ for all $t \in [0, 1, \ldots, T]$,
it follows that he will choose a  debt sequence that  satisfies 
$$ \eqalign{ b_T & \leq y_T \cr
                    b_{T-1} & \leq y_{T-1} + \beta y_T \cr
                    \vdots & \quad \quad \quad \vdots \cr
                    b_t & \leq y_t + \beta y_{t+1} + \cdots + \beta^{T-t} y_T }$$
                    
\medskip
\noindent{\bf c.} Argue that it is not necessary to impose so-called ``natural debt limits'' to make the consumer's
problem well posed.

\medskip
\noindent{\bf d.} Compute the consumer's optimal consumption plan.                                   


                                






\medskip

\noindent{\bf 1.}
Consider an infinite horizon pure exchange economy with unit measures of two types of consumers. The consumption good is not storable.  There is no uncertainty.
Consumers of type  $i=1,2$ rank consumption streams according to the utility functional
$$ \sum_{t=0}^\infty \beta^t \ln (c^i_t) $$
where $c_t^i$ is consumption of a type $i$ person at time $t \geq 0$ and $\beta \in (0,1)$. 
Each consumer of type $1$ is endowed with $(1-\lambda^t)$ units of the time $t$ consumption good where
$\lambda \in (0,1)$. Each person of type $2$ is endowed with $\lambda^t$ units of the time $t$ consumption good. 
There is a market in one-period risk-free loans each period.  Consumer $i$'s budget constraint at time $t \geq 0$ is
$$ c_t^i + b_t^i \leq R_t^{-1} b^i_{t+1} + y_t^i $$
where $b_t^i$ is one period {\it debt} that consumer $i$ owes at the beginning of period $t$ and $R_t$ is the 
equilibrium gross interest rate on one-period loans between periods $t$ and $t+1$.  Assume that $b_0^i=0$ for
$i=0,1$.

\medskip

\noindent{\bf a.} Define a notion of a ``natural borrowing limit'' for each type of consumer and compute it
for $i =1,2$.

\medskip
\noindent {\bf b.}. Define a competitive equilibrium, assuming that natural borrowing limits apply to each type of consumer.  

\medskip
\noindent{\bf c.} Compute a competitive equilibrium, being sure to include all of the objects that comprise a competitive
equilibrium.

\medskip
\noindent {\bf d.} Is a competitive equilibrium allocation Pareto optimal?

\medskip
\noindent{\bf e.} Consider an economy that is identical with the preceding one except that now 
consumers of both type are subject to the {\it ad hoc\/} borrowing limit
$$ b^i_{t+1} \leq 0 \quad \forall t \geq 0 .$$   Compute a competitive equilibrium under the {\it ad hoc\/} borrowing
limits.  Argue that the equilibrium one-period gross interest rate satisfies
$$ R_t^{-1} = \beta \max_i \left\{ {\frac{u'(c^i_{t+1})}{u'(c^i_t)}}   \right\} .$$

\medskip
\noindent{\bf f.} A friend guesses that the competitive equilibrium interest rate at which lenders are willing 
to lend is $R_t = \beta^{-1} \lambda $ for all $t \geq 0$.  She also says that there exist unsatisfied prospective borrowers
who would be willing to borrow at gross risk-free interest rate $\tilde R_t = \beta^{-1} \left({\frac{1 - \lambda^{t+1}}
{1 - \lambda^t}}\right) > R_t$.  Do you agree? Please explain why or why not.

\medskip
\noindent{\bf g.} Now consider a ``small open economy'' that is identical with the one described in part {\bf e} except
that now lenders have the opportunity to make non-negative  one-period risk-free loans at the ``world interest rate''
$\hat R_t = \beta^{-1}$ for all $t \geq 1$.  Please compute a competitive equilibrium for this economy.  

\medskip
\noindent{\bf h.} Please compare competitive equilibrium allocations for economies {\bf e} and {\bf g} according to a Pareto
criterion.


\medskip

\noindent{\bf 2.} This exercise is designed to  indicate why it is {\it not\/} necessary to impose ``natural borrowing limits'' 
in  finite-horizon versions of the pure-exchange economies that we have studied in this chapter. Consider a pure exchange economy with sequential trading of one-period Arrow securities. Let $u(c)$ be a one-period utility 
function satisfying $u' > 0, u'' < 0$ and $\lim_{c \downarrow 0} u'(c) = 0$.
Consumer $i$ has preferences ordered by
$$ \sum_{t=0}^T \sum_{s_t} \beta^t u(c^i_t(s_t)) \pi_t(s^t)  $$
where $s_t \in {\bf S}$ is governed by a time-invariant Markov chain with one step-transition density
$\pi(s' | s)$ and $\pi_t(s^t) = \pi(s_t| s_{t-1}) \pi(s_{t-1}| s_{t-2}) \cdots \pi(s_1 | s_0) \pi_0(s_0)$ and
consumer $i$'s endowment is a time invariant function of the Markov state so that
$y_t^i = y^i(s_t)$. The aggregate endowment $y_t = \sum_i y^i(s_t) = y(s_t)$.   We assume that consumption
claims at date $t$ conditional on Markov state $s_t$ are traded one-period in advance (i.e., there are complete markets in
one-period Arrow securities that trade at competitive equilibrium prices $Q_t(s_{t+1} | s_t), t =0, 1, \ldots, T-1$, it being
understood that $Q_T(s_{T+1}|s_T) = 0$ identically). 


Let's study the problem of  a consumer of type $i$, temporarily suppressing the $i$ superscripts everywhere 
to keep things tidy.  The consumer has list of optimal value functions $v_t(a, s)$ for $t = 0, 1, \ldots, T$ that satisfy
the Bellman equations
$$ v_t(a_t, s_t) = \max \left\{ u(c_t) + \beta \sum_{s_{t+1}} v_{t+1} (a_{t+1}(s_{t+1}, s_{t+1})) \pi(s_{t+1}|s_t)  \right\}$$
where  maximization is over $c_t, \{a_{t+1}(s_{t+1})\}_{s_{t+1} \in {\bf S}}$ and subject to the time
$t$ budget constraint
$$ c_t + \sum_{s_{t+1}} Q_t(s_{t+1} | s_t) a_{t+1}(s_{t+1}) \leq y(s_t) + a_t  $$
for $t < T$ and
$$ c_T \leq y(s_T) + a_T $$
for $t=T$.  The time $T$ budget constraint incorporates the fact that because the consumer will not exist
at times $t > T$, he can sell no claims on time $T+1$ consumption and would not want to buy them either.

Turn first (or,  because we are working backwards, is it ``last''?) to the Bellman equation at $T$. Solving the budget constraint at equality for $c_T =y(s_T) + a_T(s_T)$ and substituting into the right side of the Bellman equation gives
$$ v_T(a_T, s_T) = u(y(s_T) + a_T(s_T) ). $$
Staring at this equation and invoking the Inada condition that assures that $c_T(s_T) >0$  implies that the consumer
will optimally devise a plan that satisfies
$$ - a_T(s_T) \leq y_T(s_T) .$$
If we move back one period to $T-1$ and deploy  a similar argument about a Inada condition at $c_{T-1} = 0$  we can deduce from the consumer's budget constraint at $T-1$ that
$$ 0 < c_{T-1} \leq a_{T-1}(s_{T-1}) + y_{T-1}(s_{T-1}) + \sum_{s_T} Q_{T-1}(s_T | s_{T-1}) y(s_T) $$
which implies
$$ - a_{T-1}(s_{T-1})  \leq  y_{T-1}(s_{T-1}) + \sum_{s_T} Q_{T-1}(s_T | s_{T-1}) y(s_T) .$$
We can continue working backwards in this way to argue that the consumer would never {\it want\/}
to issue one-period arrow securities that exceed the continuation value of his endowment evaluated at 
prices implied  by  appropriately iterating the Arrow securities pricing kernels $Q_t(s_{t+1}|s_t)$.

Thus, with a finite horizon $T$, two conditions  allow us to dispense with so-called ``natural borrowing limits'':
({\bf i})  that it is feasible for the consumer to set $c_t >0$ for all states and dates $t \in [0, 1, \ldots, T]$;
and ({\bf ii}) that the consumer cannot {\it sell\/} claims on time $T+1$ consumption, so that $a_{T+1}(s_{T+1}) = 0$.


\medskip
\noindent {\bf 3.} Consider an infinite-horizon pure exchange economy that exists at $t=0, 1, \ldots $.
A representative consumer receives an endowment $y_t(s_t)$ of a nonstorable consumption good at time $t$
in Markov state $s_t \in {\bf S}$.  Complete markets in time- and  history-contingent consumption goods
meet at time $0$. Let $u(c)$ be a one-period utility 
function satisfying $u' > 0, u'' < 0$ and $\lim_{c \downarrow 0} u'(c) = 0$. The representative consumer ranks consumption processes according to 
$$ \sum_{t=0}^T \sum_{s^t} \beta^t u(c_t(s^t)) \pi_t(s^t)$$
where $\pi_t(s^t)$ is a joint probability function over $s^t$ induced by the Markov chain $(P, \pi_0)$ on $s_t$.
The representative consumer's budget constraint is
$$ \sum_{t=0}^\infty \sum_{s^t} q_t^0(s^t) c_t(s^t) \leq  \sum_{t=0}^\infty \sum_{s^t} q_t^0(s^t) y(s_t)  $$
 
\medskip
\noindent{\bf a.} Please define a competitive equilibrium with all trades occurring at time $0$.
\medskip
\noindent{\bf b.} Please compute a competitive equilibrium with time $0$ trades being careful to compute
all objects that comprise a competitive equilibrium.
\medskip
\noindent{\bf c.} Please define an equilibrium with sequential trades of a complete set of one-period Arrow  securities.
\medskip
\noindent{\bf d.} Please compute an equilibrium with sequential trading of a complete set of one-period Arrow securities.


\medskip
\noindent {\bf 4.} Consider an infinite-horizon pure exchange economy that exists at $t=0, 1, \ldots $.
There are unit measures of  two types of consumers.
A type 1 consumer receives an endowment $y^1(s_t) = y(s_t)$ of a nonstorable consumption good at times $t\geq T+1$
in Markov state $s_t \in {\bf S}$ and zero endowment for all $ t \leq T$.  A type 2 consumer receives an endowment $y^2(s_t) = y(s_t)$ of a nonstorable consumption good at times $t =0, \ldots, T$
in Markov state $s_t \in {\bf S}$ and zero endowment for all $t \geq T+1$.  Notice that $y^1(s_t) + y^2(s_t) = y(s_t)$ for all $t \geq 0$.

Complete markets in time- and  history-contingent consumption goods
meet at time $0$. A consumer of type 1  ranks consumption plans according to 
$$ \sum_{t=0}^T \sum_{s^t} \beta^t u(c_t^1(s^t)) \pi_t(s^t)$$
where $\pi_t(s^t)$ is a joint probability function over $s^t$ induced by the Markov chain $(P, \pi_0)$ for $s_t \in 
{\bf S}$.
A consumer of type 2 ranks consumption plans according to 
$$ \sum_{t=T+1}^\infty \sum_{s^t} \beta^t u(c_t^2(s^t)) \pi_t(s^t) .$$
%where $\pi_t(s^t)$ is a joint probability function over $s^t$ induced by the Markov chain $(P, \pi_0)$ on $s_t$.
A type $i$ consumer's budget constraint is
$$ \sum_{t=0}^\infty \sum_{s^t} q_t^0(s^t) c^i_t(s^t) \leq  \sum_{t=0}^\infty \sum_{s^t} q_t^0(s^t) y^i(s_t) . $$
 
\medskip
\noindent{\bf a.} Please define a competitive equilibrium with all trades occurring at time $0$.
\medskip
\noindent{\bf b.} Please compute a competitive equilibrium with time $0$ trades, being careful to compute
all objects that comprise a competitive equilibrium.
\medskip
\noindent{\bf c.} Please define an equilibrium with sequential trades of a complete set of one-period Arrow  securities.
\medskip
\noindent{\bf d.} Please compute an equilibrium with sequential trading of a complete set of one-period Arrow securities.



\medskip
\noindent {\bf 5.} Consider an infinite-horizon pure exchange economy that exists at $t=0, 1, \ldots $.
There are unit measures of  two types of consumers.
A type 1 consumer receives an endowment $y^1(s_t) = y(s_t)$ of a nonstorable consumption good at times $t\geq T+1$
in Markov state $s_t \in {\bf S}$.  A type 2 consumer receives an endowment $y^2(s_t) = y(s_t)$ of a nonstorable consumption good at times $t =0, \ldots, T$
in Markov state $s_t \in {\bf S}$. Please notice that $y^1(s_t) + y^2(s_t) = y(s_t)$.
Let $u(c)$ be a one-period utility 
function satisfying $u' > 0, u'' < 0$ and $\lim_{c \downarrow 0} u'(c) = 0$.
 A consumer of type 1  ranks consumption plans according to 
$$ \sum_{t=0}^T \sum_{s^t} \beta^t u(c_t^1(s^t)) \pi^1_t(s^t)$$
where $\pi^1_t(s^t)$ is a joint probability function over $s^t$ induced by the Markov chain $(P^1, \pi^1_0)$ for $s_t
\in {\bf S}$.
A consumer of type 2 ranks consumption plans according to 
$$ \sum_{t=T+1}^\infty \sum_{s^t} \beta^t u(c_t^2(s^t)) \pi^2_t(s^t) .$$
where $\pi^2_t(s^t)$ is a joint probability function over $s^t$ induced by the Markov chain $(P^2, \pi^2_0)$ for $s_t
\in {\bf S}$.
Assume that $\pi^1_t(s^t)$ and $\pi^2_t(s^t)$ put positive probabilities on the same histories $s^t$ for all $t\geq 0$.

\medskip
\noindent{\bf a.} Please pose and solve a planning problem that puts Pareto weigh $\lambda \in (0,1)$ on consumers
of type 
$1$ and $ 1 - \lambda$ on consumers of type $2$.
 
\medskip
\noindent{\bf b.} Complete markets in time- and  history-contingent consumption goods
meet at time $0$.
A type $i$ consumer's budget constraint is
$$ \sum_{t=0}^\infty \sum_{s^t} q_t^0(s^t) c^i_t(s^t) \leq  \sum_{t=0}^\infty \sum_{s^t} q_t^0(s^t) y^i(s_t) . $$Please define a competitive equilibrium with all trades occurring at time $0$.
\medskip
\noindent{\bf c.} Please compute a competitive equilibrium with time $0$ trades, being careful to compute
all objects that comprise a competitive equilibrium.
%\medskip
%\noindent{\bf c.} Please define an equilibrium with sequential trades of a complete set of one-period Arrow  securities.
%\noindent{\bf d.} Please compute an equilibrium with sequential trading of a complete set of one-period Arrow securities.

\medskip
\noindent {\bf 6.} \quad {\bf Redoing chapter 9}
\medskip
\noindent To streamline things, we assume that the aggregate endowment of a nonstorable consumption good is a time invariant function $y_t = y(s_t)$
of a Markov state $s_t \in {\bf S}$ governed by a Markov chain $P, \pi_0$.  We assume that a representative
consumer cares only about a consumption profiles $\{c_t(s_t)\}_{t=0}^\infty$ that are functions of the Markov state.
The consumer ranks consumption profiles according to the utility functional
$$ \sum_{t=0}^\infty \sum_{s_t \in {\bf S}} \beta^t u(c_t(s_t)) \pi_t(s_t | s_0) $$
where $u(c) = \left({\frac{c^{1-\gamma}}{1 - \gamma}}\right)$ with $\gamma >0$ and $\pi_t(s_t|s_0)$ is the conditional distribution over $s_t \in {\bf S}$ induced by Markov chain  $P, \pi_0$.
We assume complete markets at time $0$, after $s_0$ has been realized, in date- and Markov state-contingent
claims to consumption at dates $t \geq 0$.  The representative consumer faces the single budget constraint
$$ \sum_{t=0}^\infty \sum_{s_t \in {\bf S}} q_t^0 (s_t) c_t(s_t) \leq  \sum_{t=0}^\infty \sum_{s_t \in {\bf S}}
q_t^0 (s_t) y(s_t)  .
$$

\medskip
\noindent{\bf a.} Confirm that an equilibrium price system satisfies
$$ q_t^0(s_t) = \beta^t u'(y(s_t)) \pi_t(s_t | s_0) .$$

\medskip
\noindent{\bf b.} Verify that
$$ {\frac{q^0_{t+1}(s_{t+1})}{q^0_t(s_t)}} = \beta {\frac{u'(y(s_{t+1}))}{u'(y(s_t))}} \pi(s_{t+1}| s_t) $$
where $\pi(s_{t+1}| s_t)$ is the probability distribution of $s_{t+1}$ conditional on $s_t$.
\medskip
\noindent {\bf Hint:} Recall that laws of probability and the Markov property imply $\pi(s_{t+1} | s_0) = \pi(s_{t+1}| s_t) \pi(s_t | s_0)$. 
Write out what you are trying to compute, use this fact, and divide.

\medskip
\noindent{\bf c.}  Define  
$$  Q(s_{t+1} | s_t) =  {\frac{q^0_{t+1}(s_{t+1})}{q^0_t(s_t)}}. $$
Please interpret $ Q(s_{t+1} | s_t)$.

\medskip
\noindent{\bf d.} Please construct a heterogeneous agent economy that has the same equilibrium prices as
the preceding representative agent economy. For this new economy, please tell how to construct a competitive equilibrium allocation.




\medskip
\noindent{\bf 7.} A  pure exchange  economy consists of two
households with identical preferences  but different endowments.
A household of type $i$ has preferences that are ordered by
$$ E_0 \sum_{t=0}^\infty \beta^t u(c_{it})  , \quad
\beta \in (0,1) \leqno(1)$$
where $c_{it}$ is time $t$ consumption
of a single consumption good,  $u(c_{it}) = u_1 c_{it} -
 .5 u_2 c_{it}^2$, where
$u_1, u_2 >0$,  and $E_0$ denotes the mathematical
expectation conditioned on time $0$ information.  A household of type
$1$
has a stochastic  endowment  $y_{1t}$ of the good
governed by
$$ y_{1t+1} = y_{1t} + \sigma \epsilon_{t+1} \leqno(2) $$
where $\sigma >0$ and $\epsilon_{t+1} $ is an i.i.d.\ process
Gaussian process with mean $0$ and variance $1$.
A household of type 2 has endowment
$$ y_{2t+1} = y_{2t} - \sigma \epsilon_{t+1} \leqno(3) $$
where $\epsilon_{t+1}$ is the {\it same\/} random process as in
(2).
At time $t$, $y_{it}$ is realized before consumption at $t$ is chosen.
Assume that at time $0$, $y_{10}  = y_{20}$ and that $y_{10}$ is substantially
less than the bliss point $u_1/u_2$.  To make things computation easier,
please assume that there is no disposal of resources.

\medskip

\noindent Assume that markets are incomplete.
There is only one traded asset: a one-period risk-free
bond that both households can either purchase or issue.
The gross rate of  return on the asset between date $t$ and
date $t+1$ is $R_t$.  Household
$i$'s budget constraint at time $t$ is
$$  c_{it} + R_t^{-1} b_{it+1}  = y_{it} + b_{it} \leqno(4) $$
where $b_{it} $ is the  value in terms of time $t$ consumption
goods of household's $i$ holdings of one-period risk-free bonds.
We require that
a consumers's holdings of bonds are subject to the restriction
$$ \lim_{t\rightarrow +\infty}  \beta^t  u'(c_{it}) E b_{it+1} = 0. \leqno(5) $$
Assume that $b_{10} = b_{20} = 0$. An  incomplete markets
competitive equilibrium is a gross interest rate sequence
$\{R_t\}$, sequences of bond holdings $\{b_{it}\}$
 for $i=1,2$,  and
feasible allocations $\{c_{it}\}, i=1,2$ such that given $\{R_t\}$,
household
$i=1,2$ is  maximizing (1) subject to the sequence of budget constraints
(4) and the given initial levels of $b_{10},b_{20}$.


\medskip
\noindent{\bf a.}  Verify that there exists an incomplete markets  equilibrium in which each household simply consumes its endowment always and in which $R_t = \beta^{-1}$  and $b_{it} = 0$ for all $t \geq 0$.

\medskip
\noindent {\bf b.} Within the incomplete markets equilibrium described above, define a one-period (personal) stochastic discount factor  for time $t+1$ random payoffs for a type $i$ person as
$$ m_{i,t+1} = \beta \left( {\frac{u'(y_{i,t+1})}{u'(y_{i,t})}} \right)  $$
and define the price that a consumer of type $i$ is willing to pay at time $t$ for a sure claim to one unit of  consumption at time $t+1$ is
$$ q^i_t = E_t m_{i,t+1} $$
where $E_t$ is the mathematical expectation conditional on time $t$ information. Please justify and interpret this formula. How does $q^1_t$ compare  with
$q^2_t$?

\medskip
\noindent{\bf c.} Please conduct the following hypothetical experiment at the equilibrium allocation of the above economy.  At time
$t$, an entrepreneur thinks about marketing a security that pays one unit of time  $t+1$ consumption
provided that $\epsilon_{t+1} \geq 0$.  She wants to know how much people inside our economy would be willing to pay for this new asset if it were made available.  At the equilibrium allocation described above,
how much would consumers of types $1$ and $2$ be willing to pay for this new  asset?
Please compare how much they would be willing to pay.

\medskip
\noindent{\bf d.} Please criticize the following claim: ``Diversity in personal stochastic discount factors is a tell-tale sign
that securities markets are incomplete.''  Please tell whether you agree or disagree and explain why.

\medskip

\noindent 8. An economy consists of overlapping generations of two-period live people. At dates $t \geq 1$ there are born
$N_t >0$ identical people who are endowed with $w_1 >0$ units of a nonstorable consumption good when they
are young and $w_2 >0 $ units of a nonstorable consumption good when they are old.  In addition, at time $1$ there
are $N_0> 0$ old people each of whom is endowed with ${M \over N_0}$ units of an intrinsically worthless unbacked (fiat) currency called
dollars (it is not in anyone's utility function and is useless for production).  
Population evolves as $N_{t+1} = n N_t$ for $t \geq 0$ where $n >0$.  The consumption profile
of a person born at time $t$ is $c_t^t$ when young and $c_{t+1}^t$ when old; he/she orders profiles according to the
utility functional
$ \ln (c_t^t) + \ln (c^t_{t+1}) .$

\medskip
\noindent 
{\bf a.} Define a competitive equilibrium in which no value is attached to fiat currency, being careful to define all
of the objects that comprise a competitive equilibrium and telling who trades what with whom.

\medskip
\noindent {\bf b.} Please compute all competitive equilibria in which no value is attached to fiat currency.

\medskip
\noindent{\bf c.} Please define a competitive equilibrium in which value is attached to fiat currency, 
being careful to define all
of the objects that comprise a competitive equilibrium and telling who trades what with whom.

\medskip
\noindent{\bf d.} A ``friend'' makes the following guess.  If $n w_1 > w_2$ there exist competitive equilibria
with valued fiat currency, while if $n w_1 < w_2$ there do not.  Please either confirm or dispose of this guess.
If the guess is true, please define and compute a {\it stationary\/} equilibrium with valued fiat currency and give formulas
for the equilibrium price level and rate of return on fiat currency.

\medskip
\noindent{\bf e.} Please argue that whenever a stationary equilibrium with valued fiat currency exists, then there
also exists a continuum of equilibria with valued fiat currency, in each of which the rate of return on currency
varies over time.  Please provide explicit formulas for an equilibrium  price level as a function of time.  
{\bf Hint:} It might help to deploy a linear difference equation.

\medskip
\noindent{\bf f.} Assuming that a nonstationary equilibrium with valued fiat currency exists, please compute the
the limiting value of the rate of return on currency as $t \rightarrow + \infty$. 


\medskip

\noindent 9. An economy consists of overlapping generations of two-period lived people.  Let $T > 2$ be a finite 
integer. At dates $1 \leq  t \leq T$ there are born
$N$ identical people who are endowed with $w_1 >0$ units of a nonstorable consumption good when they
are young and $w_2>0 $ units of a nonstorable consumption good when they are old.  
At dates $t \geq T+1$ there are born
$N$ identical people who are endowed with $\bar w_1 >0$ units of a nonstorable consumption good when they
are young and $\bar w_2>0$ units of a nonstorable consumption good when they are old. In addition, at time $1$ there
are $N> 0$ old people each of whom is endowed with ${M \over N}$ units of an intrinsically worthless unbacked (fiat) currency called
dollars (it is not in anyone's utility function and is useless for production). The consumption profile
of a person born at time $t$ is $c_t^t$ when young and $c_{t+1}^t$ when old; he/she orders profiles according to the
utility functional
$ \ln (c_t^t) + \ln (c^t_{t+1}) .$


\medskip
\noindent 
{\bf a.} Define a competitive equilibrium in which no value is attached to fiat currency, being careful to define all
of the objects that comprise a competitive equilibrium and telling who trades what with whom.

\medskip
\noindent {\bf b.} Please compute all competitive equilibria in which no value is attached to fiat currency.

\medskip
\noindent{\bf c.} Please define a competitive equilibrium in which value is attached to fiat currency, 
being careful to define all
of the objects that comprise a competitive equilibrium and telling who trades what with whom.

\medskip
\noindent {\bf d.} Please state conditions on $w_1, w_2, \bar w_1, \bar w_2$ for which there exist competitive
equilibria with valued fiat currency.

\medskip
\noindent {\bf e.} Under sufficient conditions for there to exist competitive equilibria with valued fiat currency,
please find explicit formulas for the equilibrium price level sequence in the equilibrium for which the value
of currency is highest for all $t \geq 1$. {\bf Hint:} It might help to deploy a linear difference equation.

\medskip
\noindent 
{\bf 10.} An economy consists of overlapping generations of two-period live people. At dates $t \geq 0$ there are born
$N_t >0$ identical people who are endowed with $w_1 >0$ units of a nonstorable consumption good when they
are young and $w_2 >0 $ units of a nonstorable consumption good when they are old.  Population evolves as $N_{t+1} = n N_t$ for $t \geq 0$ where $n >0$ and $N_0 >0$.  The consumption profile
of a person born at time $t$ is $c_t^t$ when young and $c_{t+1}^t$ when old.

\medskip
\noindent{\bf a.} Please show that for $t \geq 1$, the frontier of feasible allocations is described by
$$ c_t^{t-1} \le n(w_1- c_t^t) + w_2, \ c_t^{t-1} \geq 0, \ c_t^t \geq 0 .$$

\medskip
\noindent{\bf b.}  For $t \geq 1$, each young person  orders consumption profiles according to the
utility functional
$ \ln (c_t^t) + \ln (c^t_{t+1}) .$
Please derive the person's offer curve.

\medskip
\noindent {\bf c.} After defining a competitive equilibrium for this economy, please draw David Gale's diagram and use it to find a Pareto optimal competitive equilibrium  allocation.

\medskip
\noindent{\bf d.} Tell how to use the same diagram to find a competitive equilibrium allocation that is not Pareto optimal.





 

\bye



